% Proposed insertions for the existing frontiers.tex; NOT a complete revision.
% No new scientific result is claimed by supplying this file.
% Confirm the canonical artifacts and integrate actual new-analysis results first.
% Replace every REVISION-PENDING marker before submission.

% Suggested replacement title:
\title[Proxy-Defined Dermatology Audit]{Auditing Proxy-Defined Subgroup Performance in a Mixed-Source Dermatology Dataset: Label Provenance, Group Support, and Fusion Baselines}

% Suggested scope/contribution paragraph (Introduction):
We investigate how subgroup performance should be interpreted when audit
attributes are inferred from the same images used for disease classification.
The study is an exploratory case study in a provider-described mixed
real/synthetic dermatology dataset. Its central contribution is an audit that
connects the provenance and meaning of proxy labels, disease composition, and
the coverage of support-dependent subgroup summaries. A comparison of simple
and metadata-conditioned fusion is a secondary analysis under the same internal
evaluation protocol. The study does not identify the effect of synthetic
augmentation, establish verified demographic fairness, or demonstrate clinical
deployment performance.

% Suggested distinction between evaluation protocols (Methods):
The archived full-data experiments used the supplied training, validation and
test folders, containing 2,000, 250 and 500 images, respectively. The reported
0.9260 accuracy for the ungated dual-stream visual head corresponds to that
500-image test set. These experiments are distinct from the seven-fold
cross-validation of the 565-image proxy-labeled subset. For the latter, disease
labels were stratified using a shuffled seven-fold splitter with seed 42.
Within each remainder, a stratified validation partition used a fraction of
0.175 and seed $42+k$ for fold $k$. This gives training/validation/test counts
of 399/85/81 in five folds and 400/85/80 in two folds.
\textbf{[REVISION-PENDING: state how identity-level matching of the compared
folds was verified, or explicitly acknowledge that only reconstruction was
possible.]}

% Suggested proxy interpretation paragraph (Methods/Limitations):
The stored age, gender-category and appearance-category outputs are
image-derived proxies. They are not verified age records, gender identity,
self-identified race or ethnicity, or clinical skin phototype. The appearance
categories retain the provenance of the automated and manual labeling process;
no mapping to Fitzpatrick types is assumed. Model-input categories are
distinguished from audit categories so that correcting an audit definition does
not silently change the inputs of an existing checkpoint. Missing, uncertain
and unmapped values are reported explicitly.

% Suggested metric definition (use only with the NEW recalculated results):
For proxy group $g$ and disease class $c$, we compute one-versus-rest
true-positive and false-positive rates as
\begin{equation}
\mathrm{TPR}_{gc}=\frac{\mathrm{TP}_{gc}}{\mathrm{TP}_{gc}+\mathrm{FN}_{gc}},
\qquad
\mathrm{FPR}_{gc}=\frac{\mathrm{FP}_{gc}}{\mathrm{FP}_{gc}+\mathrm{TN}_{gc}}.
\end{equation}
A rate with a zero denominator is undefined. Each between-group range requires
at least two groups with an estimable rate. For class $c$, our declared
one-versus-rest equalized-odds difference is
\begin{equation}
\mathrm{EOD}_{c}=\max\left\{
\max_g\mathrm{TPR}_{gc}-\min_g\mathrm{TPR}_{gc},
\max_g\mathrm{FPR}_{gc}-\min_g\mathrm{FPR}_{gc}\right\}.
\end{equation}
We report class-specific values and estimable group coverage. The unweighted
five-class macro-average is reported only when all five class values are
estimable; an average over fewer classes is explicitly labeled as such. This
definition replaces the earlier TPR-only summary and is applied consistently
across compared models.

% Suggested support/uncertainty paragraph (Methods):
We report subgroup correct counts, sample sizes and marginal 95\% Wilson
intervals. Support thresholds of 5, 10, 15, 20, 30 and 50 images are sensitivity
analyses, not validity cutoffs. For every threshold, we report the number of
retained groups and the fraction of evaluated images represented. Pooled
out-of-fold filtering is reported separately from filtering inside individual
test folds. The latter can exclude a group even when its pooled count is
substantial. Worst-group accuracy is the minimum over the explicitly retained
set. Conservative simultaneous binomial bounds are used for the minimum and
error-rate gaps, conditional on the observed predictions. These summaries do
not quantify uncertainty from model selection, annotation error or unknown
patient clustering.

% Suggested architecture clarification (Methods/Supplement):
The released metadata-conditioned DSAF head uses DINOv3 ViT-B/16 features of
dimension 768. Its two gates are feature-wise vectors of dimension 384. Each
gate maps the 256-dimensional metadata representation through a
$256\!\rightarrow\!384\!\rightarrow\!384$ network with GELU and a sigmoid,
followed by $0.5+s(\sigma(\cdot)-0.5)$. The released standard configuration uses
$s=0.5$. The gates modulate 384-dimensional global and local vectors by an
elementwise weighted sum; they are not constrained to sum to one. The local
branch in the released implementation consumes all non-CLS tokens, including
register tokens. This computation is preserved when recovering the original
checkpoints. The full-data image-only dual-stream head is a different model:
it concatenates the global and local vectors and has no metadata gates.

% Proposed interpretation of verified aggregate support analysis.
% Keep only if the historical category definitions are explicitly labeled;
% replace values if audit age categories are recomputed.
In the archived concatenation results, the within-fold threshold of at least
15 images retained only one or two composite groups and represented about
49.2\% of held-out images on average. Applying the same numerical threshold to
pooled out-of-fold counts retained 11 groups and 532 of 565 images. These
statistics describe different evaluation populations. The higher worst-group
accuracy after filtering does not reflect a change in any prediction or a
fairness intervention.

% Suggested comparison conclusion; include final paired estimates separately.
The observed comparison does not demonstrate an advantage of the adaptive
head under this dataset and protocol. It does not establish equivalence or
universal superiority of simple concatenation. Interpretation is limited by
the number of evaluated cases, overlapping cross-validation training sets,
configuration-selection history, and the meaning and provenance of the
image-derived proxy features.

% Suggested explanation limitation (Discussion):
Selected attention or saliency maps are descriptive visualizations. They do
not establish causal reliance on lesions, exclude background shortcuts, or
validate subgroup fairness. Quantitative claims about lesion-specific reliance
would require verified masks, a defined perturbation protocol, and suitable
controls evaluated on genuinely held-out cases.
